\documentclass[11pt]{article}
\usepackage{fontspec}
\usepackage[french]{babel}
\usepackage{amsmath}
\usepackage{amssymb}
\usepackage{amsthm}
\usepackage{mathrsfs}
\usepackage{tikz-cd}
\usepackage{enumitem}
\usetikzlibrary{babel}

\newcommand{\printedpage}[1]{\marginpar{\footnotesize p.~#1}}
\newcommand{\msp}[1]{\marginpar{\footnotesize [#1]}}
\newcommand{\Mon}{\mathrm{Mon}}
\newcommand{\Ob}{\mathrm{Ob}}
\newcommand{\Fl}{\mathrm{Fl}}
\newcommand{\Ss}{\mathrm{Ss}}
\newcommand{\card}{\mathrm{card}}
\newcommand{\cercle}[1]{\tikz[baseline=(c.base)]\node[draw,circle,inner sep=1pt](c){#1};}

\begin{document}

\printedpage{49}
% head: A. Grothendieck : Dérivateurs, Ch. XV, Édition des Universités de Montpellier II et Paris VII
% foot: Version 5 novembre 2010

\textit{et}
\[
\text{\textit{il existe une petite partie} } TC_0 \subset \Mon, \text{ \textit{telle que} } \Mon \cap W = \widetilde{TC_0}\ . \tag*{(7.2.5)} % [sic] (7.2.5) -> numéro probablement autre (peut-être (7.2.3)), le même numéro (7.2.5) est repris plus bas pour une autre formule et précède (7.2.4)
\]

\textit{(De façon plus précise, si} $\pi \geq \pi_0$ \textit{est un cardinal tel que} $W$ \textit{soit} $\pi$\textit{-accessible, on peut prendre} $TC_0 = (\Mon \cap W)^{\pi}$ \textit{(ou une partie} petite \textit{de} $(\Mon \cap W)^{\pi}$ \textit{ayant les mêmes classes d'isomorphie).) Donc} (th.~1 de §1) \textit{le couple} $(TC = \Mon \cap W,\ F = (TC)_*)$ \textit{est un couple de Quillen clos, et il y a même une factorisation} fonctorielle \textit{de toute flèche} $f$ \textit{en}
\[
f \ = \ p\, i \ ; \qquad i \ \in \ TC \ , \qquad p \ \in \ F \ . \tag*{(7.2.4)}
\]
\textit{Si on pose}
\[
C \ = \ \Mon \ , \qquad TC \ = \ \Mon \cap W \ , \qquad F \ = \ (TC)_* \ , \qquad TF \ = \ C_* \ , \tag*{(7.2.5)}
\]
\textit{le triple} $(W, C, F)$ \textit{est un triple de Quillen clos, si et seulement si on a les conditions supplémentaires suivantes.}
\[
\left\{
\begin{array}{lll}
\text{(7.2.6)} & \cercle{g}\ \text{\textit{(sur $\mathcal{M}$ seulement)}} & \parbox[t]{0.36\textwidth}{Mon \textit{stable par cochangement de base.}} \\[4ex]
\text{(7.2.7)} & \underbrace{TF}_{\overset{\text{déf}}{=}\,(\Mon)_*} \ \subset \ W & \parbox[t]{0.36\textwidth}{i.e. \textit{si} d) \textit{toute} $f : X \longrightarrow Y$ \textit{qui fait de} $X$ \textit{un objet injectif de} $\mathcal{M}/Y$\textit{, est dans} $W$ \textit{(donc aussi dans} $W_{\mathrm{univ}}$\textit{).}}
\end{array}
\right.
\]

\msp{page 90}

\textsc{Démonstration de} 7.1. On peut donner un énoncé un peu plus délicat, en supposant d'emblée donnée, non une partie $M_0$ de $\Ob \mathcal{M}$, mais une partie
\[
\Mon_0 \ \subset \ \Mon \ , \tag*{(7.1)}
\]
satisfaisant la condition suivante

\begin{quote}
Pour tout monomorphisme $i : A \longrightarrow B$ dans $\mathcal{M}$, il existe $i' \in \Mon_0$, $i' : A' \longrightarrow B'$, et une flèche $i' \longrightarrow i$ dans $\underline{\Fl}(\mathcal{M})$,
\[
\begin{tikzcd}[row sep=large, column sep=large]
A' \arrow[r, "{i'}"] \arrow[d] & B' \arrow[d] \\
A \arrow[r, "i"] & B \ ,
\end{tikzcd}
\tag*{(7.2)}
\]
telle que la flèche correspondante
\[
C \ = \ A \amalg_{A'} B' \ \longrightarrow \ B
\]
soit un monomorphisme, et que $B' \longrightarrow B$ ne se factorise pas par $A$.
\end{quote}

C'est cela qui était visé par la condition (7.1.1) sur $M_0$, qui implique (7.2) sur $\Mon_0$ défini en termes de $M_0$ comme dit dans 1°.

\printedpage{50}
% head: A. Grothendieck : Dérivateurs, Ch. XV, Édition des Universités de Montpellier II et Paris VII
% foot: Version 5 novembre 2010

Dans ce qui suit, on suppose donnée une partie $\Mon_0 \subset \Fl(\mathcal{M})$ satisfaisant (7.2). En pratique, elle sera souvent déduite d'une partie $M_0$ de $\Ob \mathcal{M}$, génératrice pour les monomorphismes, comme explicité dans 1°, mais on n'a pas besoin du fait que $i' \longrightarrow i$ soit un monomorphisme (\textit{i.e.} (7.1.1) est plus forte que ce qu'il nous faut). P. ex. si $B' \in M_0$ et $B' \longrightarrow B$ est donné, ne se factorisant pas par $A$, on pourra prendre $A' = B' \times_B A$ et $i' : A' \longrightarrow B'$ (pour assurer $i' \in \Mon_0$, il faudra donc supposer $M_0$ stable par passage aux sous-objets), et il y a des chances

\msp{page 91}

que $C \longrightarrow B$ sera alors bel et bien un monomorphisme \ldots

1°) Il faut prouver
\[
\Mon \ \subset \ \widetilde{\Mon_0}
\]
moyennant (7.2) (et la stabilité de $\mathcal{M}$ par $\varinjlim$ filtrantes). Cela va résulter directement du

\textbf{Lemme 7.3.} \textit{Soient $\mathcal{M}$ une catégorie stable par limites inductives filtrantes, et $C_0$, $C$ des parties de} $\Fl(\mathcal{M})$\textit{. On supposera}
\begin{enumerate}[label={}]
\item[\cercle{a}] \textit{$C$ formée de monomorphismes, et}
\item[\cercle{b}] \textit{que la condition suivante (généralisant} (7.2) \textit{ci-dessus) soit vérifiée.}
\end{enumerate}

\begin{quote}
\textit{Pour toute flèche $i : A \hookrightarrow B$ dans $C$ telle que $i$ ne soit pas un isomorphisme, il existe $i' : A' \longrightarrow B'$ dans $\widetilde{C_0}$, et une flèche $i' \longrightarrow i$ dans} $\underline{\Fl}(\mathcal{M})$\textit{, telle que dans le diagramme}
\[
\begin{tikzcd}[row sep=large, column sep=large]
A' \arrow[r, "{i'}"] \arrow[d, hook] & B' \arrow[d] \arrow[rd, dotted] & \\
A \arrow[r, "\alpha"] \arrow[rr, dotted, bend right, "i"'] & D \ = \ A \amalg_{A'} B' \arrow[r, "\beta"] & B \ ,
\end{tikzcd}
\tag*{(7.3.1)}
\]
\textit{la flèche $\beta$ soit dans $C$, et que $B' \longrightarrow B$ ne se factorise pas par $A$ (ou, ce qui revient au même, que le sous-objet de $B$ défini par $D$ et $\beta$ ne soit pas égal à celui défini par $A$ et $i$).}
\end{quote}

\textit{On suppose de plus}
\begin{enumerate}[label={}]
\item[\cercle{c}] \textit{$C$ stable par} $\varinjlim$ \textit{filtrantes, et}
\item[\cercle{d}] \textit{que pour tout $B \in \Ob \mathcal{M}$, l'ensemble des sous-objets de $B$ est petit.} (Condition des plus anodines! C'est une condition \textit{sur} $\mathcal{M}$, n'impliquant pas $C$).
\end{enumerate}
\textit{Alors on a}
\[
C \ \subset \ \widetilde{C_0} \ , \qquad \text{\textit{donc} } \widetilde{C} \ \subset \ \widetilde{C_0} \ , \tag*{(7.3.2)}
\]
\textit{donc, si $C_0 \subset C$, on conclut}

\msp{page 92}
\[
\widetilde{C} \ = \ \widetilde{C_0} \ . \tag*{(7.3.3)}
\]

\printedpage{51}
% head: A. Grothendieck : Dérivateurs, Ch. XV, Édition des Universités de Montpellier II et Paris VII
% foot: Version 5 novembre 2010

(Ce lemme implique aussitôt la première assertion de 7.1, 1°.) Prouvons ce lemme, \textit{i.e.} que toute $i \in C$ est dans $\widetilde{C_0}$, en construisant par récurrence transfinie une suite strictement croissante de sous-objets de $B$,
\[
A_\alpha \ \overset{j_\alpha}{\hookrightarrow} \ B \ , \qquad \text{avec les } j_\alpha \in C \ ,
\]
et avec :

\cercle{1°} $A_0 = A$. On prouvera que $A \xrightarrow{\ i_\alpha \,=\, i_{\alpha,0}\ } A_\alpha$ dans $\widetilde{C_0}$, et il y aura un ordinal final $\alpha$ tel que $A_\alpha = B$\,; on en conclura donc que $i = i_\alpha$ est dans $\widetilde{C_0}$.

Si $i$ est un isomorphisme, on sait que $i \in \widetilde{C_0}$, il n'y a rien à prouver (la construction transfinie s'arrête au cran 0). Les deux pas de la construction, après le pas initial 1°, sont :

\cercle{2°} Soit $A_\alpha$ construit. Si $A_\alpha = B$, la construction est arrêtée. Si $A_\alpha \xrightarrow{j_\alpha} B$ n'est pas un isomorphisme, on applique (7.3.1) à $j_\alpha \in C$, pour trouver un $D \overset{j}{\hookrightarrow} B$, donc $j \in C$ et $D \supsetneqq A_\alpha$. De plus,
\[
i_{\alpha+1,\alpha} \ : \ A_\alpha \ \longrightarrow \ A_{\alpha+1} \ \overset{\text{déf}}{=} \ D \qquad \text{est dans } \widetilde{C_0} \ , \tag*{(7.3.4)}
\]
car déduit de $i' \in \widetilde{C_0}$ par cochangement de base.

\cercle{3°} Si $\alpha$ est ordinal limite, on pose
\[
A_\alpha \ = \ \varinjlim_{\beta < \alpha} A_\beta \ , \tag*{(7.3.5)}
\]

\msp{page 93}

et la flèche canonique
\[
A_\alpha \ \xrightarrow{j_\alpha} \ B
\]
est dans $C$, car les $j_\beta$ le sont et $C$ est stable par limites inductives filtrantes.

On a ainsi achevé la description des systèmes transfinis. Pour tout $\alpha$, on a
\[
\card \alpha \ \leq \ \card \Ss(B) \ ,
\]
puisque l'application $\beta \longmapsto A_\beta$ de l'ensemble $I_\alpha = \{\beta \mid \beta < \alpha\}$ dans $\Ss(B)$ est \textit{strictement} croissante. Donc il existe un cardinal $\alpha$ tel que $\card \alpha \leq \card \Ss(B)$, tel que $A_\alpha = B$. Mais alors $i \in \widetilde{C_0}$, comme limite inductive (ou composé) transfinie de flèches appartenant à $\widetilde{C_0}$. (Il suffit de la condition $\forall\ \beta$, $i_{\beta+1,\beta} \in \widetilde{C_0}$ pour impliquer que tous les morphismes de transition $i_{\beta',\beta}$ sont dans $\widetilde{C_0}$.)

Revenons à 7.1, 2°. Il reste à prouver que si Mon est stable par cochangement de base, alors on a
\[
\Mon \ = \ \widetilde{\Mon_0} \quad .
\]

Mais comme Mon est déjà stable par composition, et par limites inductives filtrantes quelconques (donc aussi par facteurs directs), il s'ensuit par le corollaire 2 du théorème 1, §1 (p. 11), que $\Mon = \widetilde{\Mon_0}$\ \ . L'énoncé de factorisation est le théorème 1 du §1.

\printedpage{52}
% head: A. Grothendieck : Dérivateurs, Ch. XV, Édition des Universités de Montpellier II et Paris VII
% foot: Version 5 novembre 2010

\msp{page 94}

Prouvons 7.1, 3°, \textit{i.e.} la formule (7.1.6),
\[
W \cap \Mon \ \subset \ \widetilde{(W \cap \Mon)^{\pi}} \ .
\]

Pour ceci, on applique encore le lemme 7.3 à
\[
C_0 \ = \ (W \cap \Mon)^{\pi} \ , \qquad C \ = \ W \cap \Mon \ .
\]

Les conditions (\textit{a}), (\textit{b}), (\textit{c}), (\textit{d}) du lemme sont vérifiées. C'est clair pour toutes, sauf pour (\textit{b}), qui demande une vérification non tautologique. Soit donc
\[
f \ : \ A \hookrightarrow B \ , \qquad f \in \Mon \cap W,\ f \text{ non isomorphisme} \ .
\]

On considère $B$ comme limite inductive\footnote[33]{\textit{D'un système inductif $\pi$-adapté.}}
\[
B \ = \ \varinjlim_{I} B_i \ , \qquad I \ = \ \Ss_{\pi}(B) \ .
\]
Pour tout $i$, on pose
\[
A_i \ = \ A \cap B_i \ = \ A \times_B B_i \ ,
\]
et
\[
f_i \ : \ A_i \ \hookrightarrow \ B_i \ ,
\]
donc
\[
f \ = \ \varinjlim_{i \in I} f_i \ .
\]
On a besoin de savoir que le système inductif des flèches $f_i$ est $\pi$-adapté, ce qui nous amène à postuler la condition d'exactitude supplémentaire suivante sur $\mathcal{M}$ :

\msp{page 95}
\[
e')\quad \text{les } \varinjlim \text{ filtrantes dans } \mathcal{M} \text{ commutent au changement de base}\ , \tag*{(7.4)}
\]

il suffit même de le supposer dans le cas très particulier qu'on a envisagé ici.

Comme $W$ est $\pi$-accessible, et $(f_i)_{i \in I}$ est $\pi$-adapté, comme $\varinjlim f_i \in W$, il s'ensuit que l'ensemble $I_W$ des $i \in I$ tels que $f_i \in W$ est cofinal. Comme $A \overset{f}{\hookrightarrow} B$ n'est pas un isomorphisme, il existe donc un $i \in I_W$ tel que $B_i \hookrightarrow B$ ne se factorise pas par $A$. On a alors le diagramme
\[
\begin{tikzcd}[row sep=large, column sep=large]
A_i \arrow[r, hook, "f_i"] \arrow[d, hook] \arrow[rd, phantom, "\text{\footnotesize cocart.}"] & B_i \arrow[d, hook] \arrow[rd, dotted, hook, "\text{\footnotesize incl.}"] & \\
A \arrow[r, hook, "\alpha_i"] \arrow[rr, dotted, hook, bend right, "f"'] & C_i \arrow[r, "\beta_i"] & B \ ,
\end{tikzcd}
\]

\printedpage{53}
% head: A. Grothendieck : Dérivateurs, Ch. XV, Édition des Universités de Montpellier II et Paris VII
% foot: Version 5 novembre 2010
avec $f_i \in (C \cap W)^{\pi} = C_0$, donc $\alpha_i \in \widetilde{C_0}$. On a besoin de savoir que $\alpha_i \in W$, donc une condition supplémentaire sur $W$ :

\begin{quote}
Si on a un carré cocartésien
\[
\begin{tikzcd}[row sep=large, column sep=large]
A \arrow[r, hook, "i"] \arrow[d] & B \arrow[d] \\
A' \arrow[r, "{i'}"] & B' \ ,
\end{tikzcd}
\tag*{(7.5)}
\]
avec $i \in \Mon \cap W$, alors $i' \in W$\footnote[34]{\textbf{N.B.} Cette condition est vérifiée.}.
\end{quote}

(j'ai renforcé (7.1.4) en conséquence).

\msp{page 96}

Comme par hypothèse $f = \beta_i \alpha_i$ est dans $W$, et $\alpha_i \in W$, il s'ensuit par (7.1.4), (\textit{a}) que $\beta_i \in W$. Il faudra que nous sachions de plus que $\beta_i$ est un monomorphisme, donc $\beta_i \in \Mon \cap W = C$, ce qui achèvera d'établir la condition manquante (\textit{b}) de 7.3. C'est la condition d'exactitude (7.2.1), qu'il me faut visiblement rajouter aux hypothèses dans 7.1, 3°. Ayant fait cela, le lemme 7.3 s'applique et établit 7.1, 3°.

\textsc{Démonstration de} 7.2. Posons
\[
C \ = \ \Mon \ , \qquad TC \ = \ C \cap W \ , \qquad TC_0 \ = \ TC^{\pi} \ ,
\]
donc par 7.1, 3°, on a
\[
\underbrace{TC_0}_{\substack{\text{essentiellement}\\ \text{petite}}} \ \subset \ TC \ \subset \ \widetilde{TC_0} \ = \ \widetilde{TC} \ ,
\]
et par le théorème 1, §1 (p. 10-12), cor. 2, pour vérifier que $TC = \widetilde{TC_0}$ (\textit{i.e.} $TC$ est Q-clos à gauche), il revient au même de voir qu'il est Q-stable à gauche, \textit{i.e.} stable par composition, par cochangement de base, par $\varinjlim$ ordinales strictes, par facteurs directs. Or les deux derniers sont impliqués par le fait que $TC = C \cap W$ est stable par $\varinjlim$ filtrantes, vu que $W$ l'est, et $C$ aussi par la condition \cercle{e} (p. 76) sur $\mathcal{M}$. D'autre part, $C = \Mon$ et $W$ étant stables par composition, itou pour $C \cap W$. Donc il reste le cochangement de

\msp{page 97}

base, et c'est assuré par 7.2.3, (\textit{c}) (dont on avait déjà besoin dans 7.1, 3°, pour établir $C \cap W \subset \widetilde{TC_0}$). (\textbf{N.B.} les hypothèses de 7.2 sont pratiquement celles de 7.1, 3°, sauf qu'on suppose que $W$ est un localiseur (sous-entendu : modérément saturé), de façon précise on n'a utilisé, jusqu'à présent, comme hypothèse supplémentaire, que la stabilité de $W$ par composition.)

Il reste à établir que $(W, C = \Mon, F = (TC)_*)$ est un triple de Quillen clos, si et seulement si les conditions (7.2.6) et (7.2.7), \textit{i.e.}
\[
\left\{
\begin{aligned}
& C = \Mon \ \ \text{stable par cochangement de base,} \\
& TF = C_* = (\Mon)_* \subset W,
\end{aligned}
\right.
\]

\end{document}
